% Black-only interior profile for absorbent uncoated and groundwood paper.
% Use the gray model so the press PDF contains no process-colour separations.
% The 12 percent title tint is deliberately above fragile highlight-screen
% values; all meaning also remains present in rules, labels, and type styles.
\definecolor{BookAccent}{gray}{0}
\definecolor{BookAccentDark}{gray}{0}
\definecolor{BookTint}{gray}{0.88}
\definecolor{CodeBackground}{gray}{1}
\definecolor{CodeComment}{gray}{0}
\definecolor{CodeKeyword}{gray}{0}

% Semantic names used in the prose.
\newcommand{\pkg}[1]{\textsf{#1}}
\newcommand{\file}[1]{\texttt{#1}}
\newcommand{\code}[1]{\texttt{#1}}
\newcommand{\cmd}[1]{\texttt{\textbackslash#1}}
\newcommand{\required}[1]{\texttt{\{#1\}}}
\newcommand{\optional}[1]{\texttt{[#1]}}
\newcommand{\placeholder}[1]{\textnormal{\textlangle#1\textrangle}}
\newcommand{\term}[1]{\emph{#1}\index{#1}}
\newcommand{\menu}[1]{\textsf{#1}}
\DeclareSIUnit{\inch}{in}

% Consistent semantic labels used throughout the book.
\crefname{chapter}{Chapter}{Chapters}
\Crefname{chapter}{Chapter}{Chapters}
\crefname{section}{Section}{Sections}
\Crefname{section}{Section}{Sections}
\crefname{figure}{Figure}{Figures}
\Crefname{figure}{Figure}{Figures}
\crefname{table}{Table}{Tables}
\Crefname{table}{Table}{Tables}
\crefname{equation}{Equation}{Equations}
\Crefname{equation}{Equation}{Equations}

% Lists used in every chapter.
\newenvironment{learningoutcomes}
  {\begin{tcolorbox}[
      breakable,
      colback=white,
      colframe=black,
      colbacktitle=BookTint,
      coltitle=black,
      boxrule=0.85pt,
      titlerule=0.55pt,
      arc=0pt,
      title={Learning outcomes},
      fonttitle=\bfseries,
      before skip=0.8\baselineskip,
      after skip=0.9\baselineskip]
   \begin{itemize}[leftmargin=1.5em,itemsep=0.25em,topsep=0.2em]}
  {\end{itemize}\end{tcolorbox}}

\newlist{chapterexercises}{enumerate}{1}
\setlist[chapterexercises]{
  label=\thechapter.\arabic*,
  leftmargin=2.4em,
  itemsep=0.65em,
  topsep=0.4em
}

\newenvironment{chapterchecklist}
  {\begin{itemize}[
      label={\textcolor{BookAccent}{\ensuremath{\square}}},
      leftmargin=1.8em,
      itemsep=0.35em]}
  {\end{itemize}}

\newcommand{\expectedoutput}{\subsection*{What you should see}}
\newcommand{\sourceanatomy}{\subsection*{Source anatomy}}
\newcommand{\chapterproject}{\section{Chapter project}}
\newcommand{\chaptersummary}{\section{Summary}}
\newcommand{\chaptercheck}{\section{Checklist}}
\newcommand{\chapterexercisestitle}{\section{Exercises}}

% Theorem structures used in the optional advanced mathematics chapter.
\theoremstyle{definition}
\newtheorem{definition}{Definition}[chapter]
\theoremstyle{plain}
\newtheorem{theorem}[definition]{Theorem}
\newtheorem{lemma}[definition]{Lemma}
